% Options for packages loaded elsewhere
\PassOptionsToPackage{unicode}{hyperref}
\PassOptionsToPackage{hyphens}{url}
\PassOptionsToPackage{dvipsnames,svgnames,x11names}{xcolor}
%
\documentclass[
  letterpaper,
  DIV=11,
  numbers=noendperiod]{scrartcl}

\usepackage{amsmath,amssymb}
\usepackage{iftex}
\ifPDFTeX
  \usepackage[T1]{fontenc}
  \usepackage[utf8]{inputenc}
  \usepackage{textcomp} % provide euro and other symbols
\else % if luatex or xetex
  \usepackage{unicode-math}
  \defaultfontfeatures{Scale=MatchLowercase}
  \defaultfontfeatures[\rmfamily]{Ligatures=TeX,Scale=1}
\fi
\usepackage{lmodern}
\ifPDFTeX\else  
    % xetex/luatex font selection
\fi
% Use upquote if available, for straight quotes in verbatim environments
\IfFileExists{upquote.sty}{\usepackage{upquote}}{}
\IfFileExists{microtype.sty}{% use microtype if available
  \usepackage[]{microtype}
  \UseMicrotypeSet[protrusion]{basicmath} % disable protrusion for tt fonts
}{}
\makeatletter
\@ifundefined{KOMAClassName}{% if non-KOMA class
  \IfFileExists{parskip.sty}{%
    \usepackage{parskip}
  }{% else
    \setlength{\parindent}{0pt}
    \setlength{\parskip}{6pt plus 2pt minus 1pt}}
}{% if KOMA class
  \KOMAoptions{parskip=half}}
\makeatother
\usepackage{xcolor}
\setlength{\emergencystretch}{3em} % prevent overfull lines
\setcounter{secnumdepth}{5}
% Make \paragraph and \subparagraph free-standing
\ifx\paragraph\undefined\else
  \let\oldparagraph\paragraph
  \renewcommand{\paragraph}[1]{\oldparagraph{#1}\mbox{}}
\fi
\ifx\subparagraph\undefined\else
  \let\oldsubparagraph\subparagraph
  \renewcommand{\subparagraph}[1]{\oldsubparagraph{#1}\mbox{}}
\fi


\providecommand{\tightlist}{%
  \setlength{\itemsep}{0pt}\setlength{\parskip}{0pt}}\usepackage{longtable,booktabs,array}
\usepackage{calc} % for calculating minipage widths
% Correct order of tables after \paragraph or \subparagraph
\usepackage{etoolbox}
\makeatletter
\patchcmd\longtable{\par}{\if@noskipsec\mbox{}\fi\par}{}{}
\makeatother
% Allow footnotes in longtable head/foot
\IfFileExists{footnotehyper.sty}{\usepackage{footnotehyper}}{\usepackage{footnote}}
\makesavenoteenv{longtable}
\usepackage{graphicx}
\makeatletter
\def\maxwidth{\ifdim\Gin@nat@width>\linewidth\linewidth\else\Gin@nat@width\fi}
\def\maxheight{\ifdim\Gin@nat@height>\textheight\textheight\else\Gin@nat@height\fi}
\makeatother
% Scale images if necessary, so that they will not overflow the page
% margins by default, and it is still possible to overwrite the defaults
% using explicit options in \includegraphics[width, height, ...]{}
\setkeys{Gin}{width=\maxwidth,height=\maxheight,keepaspectratio}
% Set default figure placement to htbp
\makeatletter
\def\fps@figure{htbp}
\makeatother

\usepackage{booktabs}
\usepackage{longtable}
\usepackage{array}
\usepackage{multirow}
\usepackage{wrapfig}
\usepackage{float}
\usepackage{colortbl}
\usepackage{pdflscape}
\usepackage{tabu}
\usepackage{threeparttable}
\usepackage{threeparttablex}
\usepackage[normalem]{ulem}
\usepackage{makecell}
\usepackage{xcolor}
\KOMAoption{captions}{tableheading}
\setcounter{figure}{0} % Reset the figure counter to 0
\setcounter{table}{0} % Reset the figure counter to 0
\renewcommand{\thefigure}{S\arabic{figure}}
\renewcommand{\thetable}{S\arabic{table}}
\makeatletter
\@ifpackageloaded{caption}{}{\usepackage{caption}}
\AtBeginDocument{%
\ifdefined\contentsname
  \renewcommand*\contentsname{Table of contents}
\else
  \newcommand\contentsname{Table of contents}
\fi
\ifdefined\listfigurename
  \renewcommand*\listfigurename{List of Figures}
\else
  \newcommand\listfigurename{List of Figures}
\fi
\ifdefined\listtablename
  \renewcommand*\listtablename{List of Tables}
\else
  \newcommand\listtablename{List of Tables}
\fi
\ifdefined\figurename
  \renewcommand*\figurename{Figure}
\else
  \newcommand\figurename{Figure}
\fi
\ifdefined\tablename
  \renewcommand*\tablename{Table}
\else
  \newcommand\tablename{Table}
\fi
}
\@ifpackageloaded{float}{}{\usepackage{float}}
\floatstyle{ruled}
\@ifundefined{c@chapter}{\newfloat{codelisting}{h}{lop}}{\newfloat{codelisting}{h}{lop}[chapter]}
\floatname{codelisting}{Listing}
\newcommand*\listoflistings{\listof{codelisting}{List of Listings}}
\makeatother
\makeatletter
\makeatother
\makeatletter
\@ifpackageloaded{caption}{}{\usepackage{caption}}
\@ifpackageloaded{subcaption}{}{\usepackage{subcaption}}
\makeatother
\ifLuaTeX
  \usepackage{selnolig}  % disable illegal ligatures
\fi
\usepackage{bookmark}

\IfFileExists{xurl.sty}{\usepackage{xurl}}{} % add URL line breaks if available
\urlstyle{same} % disable monospaced font for URLs
\hypersetup{
  pdftitle={Engaging students with statistics through choice of real data context on homework},
  colorlinks=true,
  linkcolor={blue},
  filecolor={Maroon},
  citecolor={Blue},
  urlcolor={Blue},
  pdfcreator={LaTeX via pandoc}}

\title{Engaging students with statistics through choice of real data
context on homework}
\usepackage{etoolbox}
\makeatletter
\providecommand{\subtitle}[1]{% add subtitle to \maketitle
  \apptocmd{\@title}{\par {\large #1 \par}}{}{}
}
\makeatother
\subtitle{Supplemental Materials}
\author{}
\date{}

\begin{document}
\maketitle

\renewcommand*\contentsname{Table of contents}
{
\hypersetup{linkcolor=}
\setcounter{tocdepth}{3}
\tableofcontents
}
\newpage

\section{Additional results}\label{additional-results}

\begin{table}

\caption{\label{tbl-r1-results}Estimates from generalized estiamting
equation fit for research question 1's investigation}

\centering{

\centering
\begin{tabular}[t]{lllll}
\toprule
Coefficient & Estiamte & Standard Error & Wald Statistic & P-value\\
\midrule
\cellcolor{gray!10}{Intercept} & \cellcolor{gray!10}{1.019} & \cellcolor{gray!10}{1.546} & \cellcolor{gray!10}{0.434} & \cellcolor{gray!10}{0.510}\\
I (Choice) & -1.229 & 1.041 & 1.394 & 0.238\\
\cellcolor{gray!10}{I (Section B)} & \cellcolor{gray!10}{1.133} & \cellcolor{gray!10}{2.052} & \cellcolor{gray!10}{0.305} & \cellcolor{gray!10}{0.581}\\
I (Week 2) & -0.195 & 1.209 & 0.026 & 0.872\\
\cellcolor{gray!10}{I (Week 3)} & \cellcolor{gray!10}{-1.328} & \cellcolor{gray!10}{1.561} & \cellcolor{gray!10}{0.723} & \cellcolor{gray!10}{0.395}\\
\addlinespace
I (Week 4) & -1.373 & 1.054 & 1.696 & 0.193\\
\bottomrule
\end{tabular}

}

\end{table}%

\newpage

\begin{longtable}[t]{>{\raggedright\arraybackslash}p{1in}>{\raggedright\arraybackslash}p{1in}>{\raggedright\arraybackslash}p{4in}}

\caption{\label{tbl-r2-themes}Themes identified in research question 2's
invesitgation.}

\tabularnewline

\toprule
Theme & Subtheme & Subtheme Description\\
\midrule
\cellcolor{gray!10}{burden of choice} & \cellcolor{gray!10}{time} & \cellcolor{gray!10}{Student noted the extra time involved with making a choice of data context}\\
burden of choice & decision & Student felt it was difficult to choose a homework\\
\cellcolor{gray!10}{collaborating} & \cellcolor{gray!10}{peers} & \cellcolor{gray!10}{Peers were working on different assignments, so sharing work or getting help was harder}\\
no difference & unaffected & Student did not feel choice made any difference on their experience with the homework\\
\cellcolor{gray!10}{no difference} & \cellcolor{gray!10}{same} & \cellcolor{gray!10}{Student thought the choice of data context was unimportant since the investigation steps were the same}\\
\addlinespace
no difference & uninteresting & Students were not interested in any of the data contexts\\
\cellcolor{gray!10}{variety} & \cellcolor{gray!10}{variety} & \cellcolor{gray!10}{Student appreciated variety of data contexts}\\
variety & more choice & Student wanted more choice in future assignments or classes\\
\cellcolor{gray!10}{relevant to student} & \cellcolor{gray!10}{relatable} & \cellcolor{gray!10}{Student liked having data they could relate to}\\
relevant to student & academic & Students like data relevant to student's academic interests\\
\addlinespace
\cellcolor{gray!10}{relevant to student} & \cellcolor{gray!10}{interesting} & \cellcolor{gray!10}{Students like having data contexts they found interesting}\\
previous knowledge & familiar & Student factored knowledge of data context into choice\\
\cellcolor{gray!10}{better experience} & \cellcolor{gray!10}{motivating} & \cellcolor{gray!10}{Student felt more motivated to do the homework}\\
better experience & engaging & Student felt the homework was more engaging\\
\cellcolor{gray!10}{better experience} & \cellcolor{gray!10}{fun} & \cellcolor{gray!10}{Student expressed that the homework was more fun}\\
\addlinespace
better experience & easier & Student felt it was easier to work on homework\\
\cellcolor{gray!10}{better experience} & \cellcolor{gray!10}{challenge} & \cellcolor{gray!10}{Student felt the homework was challenging}\\
better experience & focus & The student felt they were able to better focus on assignment\\
\cellcolor{gray!10}{learn more about topic} & \cellcolor{gray!10}{learning} & \cellcolor{gray!10}{Student enjoyed learning about data context through assignment}\\
context is valuable & helpful & Student felt the data context helped them understand the statistical concepts better\\
\addlinespace
\cellcolor{gray!10}{statistics is valuable} & \cellcolor{gray!10}{meaningful} & \cellcolor{gray!10}{Student felt interesting data contexts helped make statistics meaningful instead of just being numbers}\\
statistics is valuable & real world & Student expressed that statistics could have real world applications\\
\cellcolor{gray!10}{choice is valuable} & \cellcolor{gray!10}{autonomy} & \cellcolor{gray!10}{Student appreciated having autonomy}\\
practice & practice & Student appreciated the choice providing more practice through non chosen assignments\\
\bottomrule

\end{longtable}

\newpage

\begin{longtable}[t]{>{\raggedright\arraybackslash}p{1in}>{\raggedright\arraybackslash}p{1in}>{\raggedright\arraybackslash}p{4in}}

\caption{\label{tbl-r3-themes}Themes identified in research question 3's
invesitgation.}

\tabularnewline

\toprule
Theme & Subtheme & Subtheme Description\\
\midrule
\cellcolor{gray!10}{current topic} & \cellcolor{gray!10}{popular} & \cellcolor{gray!10}{Topic commonly discussed/popular}\\
current topic (converse) & old & Student felt data was too old\\
\cellcolor{gray!10}{important} & \cellcolor{gray!10}{important} & \cellcolor{gray!10}{Student felt the topic was important}\\
important & ethics & Student expressed appriciation for data context being related to ethics\\
\cellcolor{gray!10}{comfort} & \cellcolor{gray!10}{political} & \cellcolor{gray!10}{Commonly discussed political/controversial topic}\\
\addlinespace
comfort & negative & Student perceives topic to have a negative connotation\\
\cellcolor{gray!10}{comfort} & \cellcolor{gray!10}{comfort} & \cellcolor{gray!10}{Student felt uncomfortable about the topic}\\
relevant & generalizability & Felt topic was relevant to large population of people\\
\cellcolor{gray!10}{relevant} & \cellcolor{gray!10}{relatable} & \cellcolor{gray!10}{Student expressed data context to be relatable with no reason}\\
relevant & major & Relevant to student's field of study or academic interests\\
\addlinespace
\cellcolor{gray!10}{relevant} & \cellcolor{gray!10}{career} & \cellcolor{gray!10}{Relevant to student's future career plans}\\
relevant & interest & Relevant to personal interests\\
\cellcolor{gray!10}{relevant} & \cellcolor{gray!10}{dislike} & \cellcolor{gray!10}{Relevant to student's dislike}\\
relevant & experience & Personal experience with topic\\
\cellcolor{gray!10}{relevant} & \cellcolor{gray!10}{student} & \cellcolor{gray!10}{Relevant to daily life of a student}\\
\addlinespace
relevant & life & Relevance to daily life\\
\cellcolor{gray!10}{previous knowledge} & \cellcolor{gray!10}{familiar} & \cellcolor{gray!10}{Familiar with / knowledgable about topic}\\
previous knowledge & sequel & Same topic student chose the previous week\\
\cellcolor{gray!10}{previous knowledge} & \cellcolor{gray!10}{predictable} & \cellcolor{gray!10}{Student felt topic/results would be predictable}\\
previous knowledge & preconception & Student has a preconception about the results from the data exploration\\
\addlinespace
\cellcolor{gray!10}{previous knowledge} & \cellcolor{gray!10}{unclear} & \cellcolor{gray!10}{Student felt unclear about what the context was}\\
learn more & intriguing & The topic was intriguing to a student\\
\cellcolor{gray!10}{learn more} & \cellcolor{gray!10}{curious} & \cellcolor{gray!10}{Curious to learn more about topic through homework}\\
learn more & results & Expressed interest in what they learned about the results\\
\cellcolor{gray!10}{learn more} & \cellcolor{gray!10}{surprise} & \cellcolor{gray!10}{Student expressed they were surprised by results}\\
\addlinespace
better experience & difficulty & Perceived topic as easy to understand\\
\cellcolor{gray!10}{better experience} & \cellcolor{gray!10}{pace} & \cellcolor{gray!10}{Student thought homework would take less time with the topic}\\
better experience & fun & Thought it would be fun to do hw on topic\\
\cellcolor{gray!10}{better experience} & \cellcolor{gray!10}{engaging} & \cellcolor{gray!10}{Thought topic would make homework more engaging}\\
better experience & motivating & Thought topic would make the homework more motivating\\
\addlinespace
\cellcolor{gray!10}{variety} & \cellcolor{gray!10}{different} & \cellcolor{gray!10}{Different from topics commonly discussed in class}\\
variety & multivariable & Student expressed appriciation for the data being multivariable\\
\cellcolor{gray!10}{random} & \cellcolor{gray!10}{random} & \cellcolor{gray!10}{Random choice}\\
random & first & First option student saw for hw\\
\cellcolor{gray!10}{data quality} & \cellcolor{gray!10}{data quality} & \cellcolor{gray!10}{Distrust of quality of data due to topic}\\
\addlinespace
data quality & self-reporting bias & Student mentioned distrust of data due to concerns of self-reporting bias\\
\cellcolor{gray!10}{data quality} & \cellcolor{gray!10}{sample size} & \cellcolor{gray!10}{Student mentioned distrust of results due to concerns of sample size}\\
\bottomrule

\end{longtable}

\section{Resources for finding real data sets for educational
purposes}\label{sec-resources}

Below is a list of resources that have been recommended to find data for
educational purposes. A dynamic version of this list of resources is
available at \href{https://catalinamedina.github.io/resources.html}.

\subsection{Repositories From Professional Organizations for
Teaching}\label{repositories-from-professional-organizations-for-teaching}

\begin{itemize}
\tightlist
\item
  \href{https://causeweb.org/jedi/}{JEDI} JEDI-informed teaching of
  statistics
\item
  \href{https://docs.google.com/spreadsheets/d/1tuCXok_y_zQpSDGssaypETf9lnB6nCh9XVoCiMaPCuc/edit?gid=591469695\#gid=591469695}{SSDSE
  Dataset Repository}
\item
  \href{https://www.causeweb.org/cause/resources/library?field_material_type_tid=104}{CAUSE
  Resource Library}
\item
  \href{https://www.causeweb.org/tshs/}{Teaching of Statistics in the
  Health Sciences (TSHS) Resources Portal}, a source for well-documented
  health-related datasets and teaching materials.
\item
  \href{http://wise.cgu.edu/helpful-links/data-sources/}{WISE (Web
  Interface for Statistics Education)} has a collection of
  demonstrations and tutorials.
\end{itemize}

\subsection{Published Research}\label{published-research}

\begin{itemize}
\tightlist
\item
  \href{https://studyfinds.org/}{Study Finds} More accessible topics
  from a wide range of fields
\item
  Google News Health and Sciences Section
\item
  \href{https://www.icpsr.umich.edu/web/pages/ICPSR/index.html}{ICPSR}
  International social research
\item
  \href{https://osf.io/}{OSF open data repository} General platform for
  sharing research, data, materials
\item
  Journals such as \href{https://journals.plos.org/plosone/}{PLOS1}
\end{itemize}

\subsection{Other Repositories or Data
Portals}\label{other-repositories-or-data-portals}

\begin{itemize}
\tightlist
\item
  \href{https://vincentarelbundock.github.io/Rdatasets/}{Rdatasets}
  Rdatasets is a collection of 3451 datasets which were originally
  distributed alongside the statistical software environment R and some
  of its add-on packages. The goal is to make these data more broadly
  accessible for teaching and statistical software development.
\item
  \href{http://www.nationmaster.com/}{Nationmaster} portal to
  international economic, demographic and social data.
\item
  \href{https://github.com/caesar0301/awesome-public-datasets}{Awesome
  Public Data Sets} by Xiaming Chen and other contributors
\item
  \href{http://dasl.datadesk.com/}{Data and Story Library (DASL)}
\item
  \href{https://github.com/awesomedata/awesome-public-datasets}{Awesome
  public datasets}
\item
  \href{https://gss.norc.org/us/en/gss/get-the-data.html}{General Social
  Survey} American demographics, behaviors, and opinions (good for two
  categorical variables)
\item
  \href{https://www.bikeshare.com/data/}{Bikeshare data portal}
\item
  \href{https://www.data.gov/}{Data.gov}
\item
  \href{https://docs.google.com/spreadsheets/d/1wZhPLMCHKJvwOkP4juclhjFgqIY8fQFMemwKL2c64vk/edit\#gid=0}{Data
  is Plural}
\item
  \href{https://edinburghopendata.info/}{Edinburgh Open Data}
\item
  \href{https://corgis-edu.github.io/corgis/}{CORGIS: The Collection of
  Really Great, Interesting, Situated Datasets}
\item
  \href{https://datasetsearch.research.google.com/}{Google Dataset
  Search}
\item
  \href{https://dataverse.harvard.edu/}{Harvard Dataverse}
\item
  \href{https://www.opendata.nhs.scot/}{NHS Scotland Open Data}
\item
  \href{https://ipums.org/}{IPUMS survey data from around the world}
\item
  \href{https://data.lacity.org/browse}{Los Angeles Open Data}
\item
  \href{https://opendata.cityofnewyork.us/}{NYC OpenData}
\item
  \href{https://statistics.gov.scot/home}{Open access to Scotland's
  official statistics}
\item
  \href{https://www.icpsr.umich.edu/icpsrweb/content/ICPSR/fenway.html}{PRISM
  Data Archive Project}
\item
  \href{https://archive.ics.uci.edu/}{UCI machine learning repository}
  Machine learning
\item
  \href{http://data.un.org/}{UN data}
\item
  \href{https://data.gov.uk/}{UK Government Data}
\item
  \href{https://sctyner.github.io/static/presentations/Misc/GraphicsGroupISU/2018-11-16-us-govt-data.html}{US
  Government Data}
\item
  \href{https://chronicdata.cdc.gov/Youth-Risk-Behaviors/DASH-Youth-Risk-Behavior-Surveillance-System-YRBSS/q6p7-56au}{Youth
  Risk Behavior Surveillance System (YRBSS)}
\item
  \href{https://rebeccabarter.com/blog/2023-03-28-data_sources}{Rebecca
  Barter's List of Public Databases}
\end{itemize}

\subsection{Open Access Textbooks}\label{open-access-textbooks}

\begin{itemize}
\tightlist
\item
  Peter K. Dunn (2024). Scientific Research and Methodology: An
  introduction to quantitative research in science and health.
  \url{https://bookdown.org/pkaldunn/SRM-Textbook}
\item
  \href{https://www.bayesrulesbook.com/}{Bayes Rules! An Introduction to
  Applied Bayesian Modeling}
\item
  \href{https://www.openintro.org/book/os/}{Open Intro Statistics}
\end{itemize}

\subsection{Other Ideas of Places to
Look}\label{other-ideas-of-places-to-look}

\begin{itemize}
\tightlist
\item
  \href{https://www.tandfonline.com/doi/abs/10.1080/00031305.1990.10475726}{Singer
  1990} suggests looking at published papers from local dissertations
\end{itemize}



\end{document}
